\externaldocument[nocite]{build/text/chapters/our-games}[url=text/chapters/our-games]
%===============================================================
\chapter{Launching and Testing the Programme}\label{chap:launch-testing}
%===============================================================

\begin{chapterabstract}
	To ensure the long-term viability of our student game incubation programme, we plan to launch it with broad outreach to effectively validate our approach with direct user feedback and adapt on the go. We launched a dedicated website to inform, register, and engage potential participants. To raise awareness, we ran targeted campaigns using posters put up at our faculty (FIT) and in the relevant Discord channels. We began user sentiment testing to understand developer priorities and expectations. For pilot testing, we are recruiting participants through direct outreach to experienced game authors, promotions at events (e.g., GameJams), and open applications at CTU. A year-long pilot phase will be followed by a structured evaluation to assess participant benefit, content relevance, delivery methods, and operational feasibility.
\end{chapterabstract}

To ensure long-term viability of our incubation programme, we aim to launch it with sizable outreach. Our strategy is to gather broad up-front feedback in our community and to start broadly inviting participants after launch. This approach will help us verify the assumptions we made when designing the mechanisms and iteratively improve.

%===============================================================
\section{User Sentiment Testing}
Understanding what student developers need, prioritise and how they perceive the programme is essential to its successful launch. Therefore, we have started gathering opinions and conducting user sentiment testing.

%---------------------------------------------------------------
\subsection{Presenting through a Website}
To direct our outreach, we have created a dedicated website [which can be found in the appendix] and was designed to directly fit into the Grafit design language from [REPO].
\begin{figure}[H]
    \includegraphics[width=0.9\textwidth]{web_landing.png}
    \caption{The grafit.games landing page}
    \label{fig:web_landing}
\end{figure}

This platform serves several key purposes:
\begin{itemize}
    \item Inform potential participants about the goals, structure, and benefits of the programme:
    \begin{figure}[H]
        \includegraphics[width=0.9\textwidth]{web_our-proposition.png}
        \caption{Screenshot from grafit.games - Our Proposition}
        \label{fig:web_our-proposition}
    \end{figure}
    \item Provide updates on the programme’s progress and outline upcoming events or changes:
    \begin{figure}[H]
        \includegraphics[width=0.9\textwidth]{web_current-events.png}
        \caption{Screenshot from grafit.games - Current}
        \label{fig:web_current-events}
    \end{figure}
    \item Allow interested individuals to register or express their intent to contribute. We are collecting a name, email (required), discord username, game description, input on what the interviewee would like to see from us, whether they would like to participate in the preparation of Grafit.games and whether they would like to be informed about updates:
    \begin{figure}[H]    
        \includegraphics[width=0.9\textwidth]{web_get-notified.png}
        \caption{Screenshot from grafit.games - Get Notified}
        \label{fig:web_get-notified}
    \end{figure}
\end{itemize}

This website serves as a communication channel, evokes credibility, and promises progress to potential participants and backers.
        
%---------------------------------------------------------------
\subsection{Campaign}
% \begin{figure}[H]
%     \noindent
%     \begin{minipage}[t]{0.6\textwidth}
%         \vspace{0pt}
%         \centering
%         \includegraphics[width=\textwidth]{posterPrint.png}
%         \caption{The grafit.games poster}
%         \label{fig:poster}
%     \end{minipage}%
%     \hfill
%     \begin{minipage}[t]{0.4\textwidth}
%         \vspace{0pt}
%         To direct interested individuals to our website, we’ve put up posters on the premises of our faculty and in the relevant discord servers. We’ve tracked the performance of our campaign using bitly.
%     \end{minipage}
% \end{figure}

\begin{figure}[H]
    To direct interested individuals to our website, we’ve put up posters on the premises of our faculty and in the relevant discord servers. We’ve tracked the performance of our campaign using bitly.
    
    \vspace{1em}
    \includegraphics[width=0.8\textwidth]{posterPrint.png}
    \caption{The grafit.games poster}
    \label{fig:poster}
\end{figure}

%===============================================================
\section{Pilot Run}
For a pilot run, we will recruit a group of developers we believe to be representative, aiming to gain a broad range of feedback. We will focus on multiple sources for our enlistments:
\begin{itemize}
    \item \textbf{Direct Outreach to Game Authors:} We will directly approach creators of exemplary games - some of which have been referenced in section~\ref{sec:example-games}. These individuals have firsthand experience with the development process within the academic environment, can provide informed feedback and crucially possess a game in the late stages of development that could be commercialized.
    \item \textbf{Engagement at Events:} We will invite competitors at hackathons such as GameJam and advertise the activities of the incubator at promotional events.
    \item \textbf{Open Applications at CTU:} Whilst we mostly expect interest from students at FIT, the programme will be available to all associates of CTU - students, staff, and alumni alike.
\end{itemize}

%===============================================================
\section{Evaluating the Support Mechanism}
We want to improve as much as possible throughout a year-long pilot run. For the consecutive summer however, a structured evaluation phase is planned. We will assess the strengths, weaknesses, impact, and sustainability of the programme and make modifications if necessary.

%---------------------------------------------------------------
\subsection{Our Goals - The Criteria We Will Evaluate}
These criteria will be included in the participant questionnaires and then consider in the evaluation:
\begin{itemize}
    \item \textbf{Participant Benefit:}  We aim to determine whether participants found the programme helpful and relevant to their needs.
    \item \textbf{Content Evaluation:} We will assess whether the teaching materials cover necessary technical, legal, and business topics without overwhelming participants or distracting from their studies.
    \item \textbf{Delivery Optimization:} We will evaluate whether the chosen mediums (e.g., workshops, online tutorials, consultations) are effective in conveying the content.
\end{itemize}

Additionally, we will determine and consider:
\begin{itemize}
    \item \textbf{Staffing Feasibility:} We will review whether we have or will be capable of attaining sufficient staffing to operate the incubator
    \item \textbf{Financial Feasibility:} We will review the expenses needed to run the incubator and whether we can attract financial backers
    \item \textbf{Feasibility of the Selected Legal Form:} We will evaluate whether the legal and operational framework poses any risk to the programme’s longevity.
\end{itemize}

%---------------------------------------------------------------
\subsection{Procedure After Evaluation}
Based on the findings from the year-long operations, we will be prepared to take decisive action:
\begin{itemize}
    \item \textbf{Adjustments in the Teaching:} In the case that 30\% of the respondents considered the programme unhelpful, or major insufficiencies in the taught information - whether it is the contents, structure, or delivery - have been identified by evaluating the questionnaires, the programme will be restructured for a consecutive run the next year.
    \item \textbf{Changing the Legal Structure:} If it is determined that the legal, executive, or operational framework does not fit the programme’s needs, it will be replaced.
    \item \textbf{Termination and Redesign:} Should the programme be deemed unhelpful by 60\% of the respondents, or should it be deemed unsustainable in its current form, it will be canceled and - hopefully - an alternative solution that better fits the needs of the community will be proposed.
\end{itemize}

We hope that such flexibility will allow us to achieve the quality and impact we envision for the initiative.
